\documentclass[conference]{IEEEtran}
\IEEEoverridecommandlockouts

\usepackage{cite}
\usepackage{amsmath,amssymb,amsfonts}
\usepackage{algorithmic}
\usepackage{graphicx}
\usepackage{textcomp}
\usepackage{xcolor}
\usepackage{booktabs}
\usepackage{array}
\usepackage[hidelinks]{hyperref}

\begin{document}

\title{Optimizing Vector Retrieval Latency and Accuracy in Resource-Constrained Explainability-Driven Offline RAG Systems}

\author{\IEEEauthorblockN{First Author}
\IEEEauthorblockA{\textit{Department of Computer Science} \\
\textit{University Name}\\
City, Country \\
first.author@university.edu}
\and
\IEEEauthorblockN{Second Author}
\IEEEauthorblockA{\textit{Department of Computer Science} \\
\textit{University Name}\\
City, Country \\
second.author@university.edu}
\and
\IEEEauthorblockN{Third Author}
\IEEEauthorblockA{\textit{Department of Computer Science} \\
\textit{University Name}\\
City, Country \\
third.author@university.edu}
}

\maketitle

\begin{abstract}
Retrieval-Augmented Generation (RAG) systems mitigate hallucination by grounding large language model (LLM) outputs in retrieved evidence, yet offline deployment on resource-constrained hardware introduces substantial retrieval latency that compounds with downstream verification costs. Existing dense-retrieval implementations using exhaustive flat-index search scale poorly with corpus size and lack the speed required for real-time explainability-driven hallucination detection. This paper presents a systematic investigation of vector indexing optimizations for an offline, explainability-augmented RAG framework in which each generated sentence must be attributed to source chunks via cosine similarity before release. We benchmark four retrieval strategies—FAISS FlatIP (baseline), Inverted File with Product Quantization (IVF-PQ), Hierarchical Navigable Small World (HNSW) graphs, and a hybrid dense-sparse approach combining FAISS with BM25—across latency, memory footprint, and Recall@K metrics on a corpus of 9,800 document chunks. Experimental results on a consumer-grade CPU demonstrate that HNSW reduces mean retrieval latency from 47.3\,ms to 3.8\,ms (12.4$\times$ speedup) while maintaining Recall@5 $\geq$ 0.93, and that hybrid retrieval achieves the highest Recall@5 of 0.96 at 8.2\,ms latency. IVF-PQ achieves 6.7$\times$ memory compression (384\,MB $\to$ 57\,MB) with acceptable Recall@5 degradation to 0.89. We analyze the Pareto frontier of the latency-accuracy-memory trade-space and provide deployment guidelines for resource-constrained explainability-driven RAG systems. Our optimizations reduce end-to-end query latency from 9.8\,s to 9.76\,s, with retrieval contributing less than 1\% of total processing time, revealing that Natural Language Inference (NLI) verification (46.5\% of latency) represents the dominant bottleneck for future optimization.
\end{abstract}

\begin{IEEEkeywords}
dense retrieval, FAISS optimization, HNSW, IVF-PQ, hybrid retrieval, retrieval-augmented generation, hallucination detection, explainability, offline NLP, resource-constrained systems
\end{IEEEkeywords}

\section{Introduction}

\subsection{Motivation and Background}

Large language models (LLMs) have demonstrated remarkable capabilities in natural language understanding and generation, yet they remain prone to \emph{hallucination}—the generation of fluent, confident, but factually incorrect or unverifiable statements \cite{ji2023survey,huang2025survey}. This phenomenon poses critical risks in safety-sensitive domains such as clinical decision support \cite{pal2023medical}, financial analysis \cite{islam2023financebench}, and legal document review, where a single fabricated claim can propagate into high-stakes downstream decisions.

Retrieval-Augmented Generation (RAG) \cite{lewis2020rag} addresses this limitation by conditioning LLM outputs on evidence retrieved from an external corpus, grounding responses in verifiable source material. While RAG substantially improves factual accuracy over purely parametric generation, benchmark studies reveal that LLMs frequently ignore, misinterpret, or extrapolate beyond retrieved evidence \cite{chen2024benchmarking,niu2024ragtruth}, necessitating explicit verification mechanisms.

Contemporary hallucination detection approaches—including sampling-based consistency checks \cite{manakul2023selfcheckgpt}, atomic fact verification \cite{min2023factscore}, and LLM-as-judge frameworks \cite{es2024ragas}—predominantly operate \emph{post hoc}, analyzing generated text after completion. Moreover, these methods commonly depend on cloud-hosted APIs or require repeated inference through large proprietary models, rendering them unsuitable for offline, air-gapped, or data-sovereign environments \cite{krainovskikh2026domain}. 

Recent work on explainable AI (XAI) for RAG systems \cite{zhang2024feel,agrawal2024attributed} has demonstrated that attribution signals—mapping generated sentences to source chunks via embedding similarity—can expose unsupported claims. However, these frameworks treat explainability as an \emph{analytical tool} rather than an \emph{architectural constraint}: attribution is computed for human inspection but does not gate model outputs. To the authors' knowledge, no existing system integrates intrinsic, real-time attribution as a mandatory hallucination control mechanism within a fully offline RAG pipeline.

\subsection{Problem Formulation}

Let $\mathcal{D} = \{d_1, \ldots, d_N\}$ denote a document corpus partitioned into $N$ overlapping chunks, each embedded as $\mathbf{e}_i \in \mathbb{R}^d$ via a sentence encoder. Given a user query $q$ with embedding $\mathbf{e}_q \in \mathbb{R}^d$, dense retrieval identifies the Top-$K$ most similar chunks:
\begin{equation}
\mathcal{R}(q) = \arg\max_{C \subset \mathcal{D}, |C|=K} \sum_{d_i \in C} \mathrm{sim}(\mathbf{e}_q, \mathbf{e}_i)
\label{eq:retrieval}
\end{equation}
where $\mathrm{sim}(\mathbf{e}_q, \mathbf{e}_i) = \frac{\mathbf{e}_q \cdot \mathbf{e}_i}{\|\mathbf{e}_q\| \|\mathbf{e}_i\|}$ is cosine similarity. A locally hosted LLM then generates a draft answer conditioned on $\mathcal{R}(q)$.

In an explainability-augmented framework, each generated sentence $s_j$ ($j=1,\ldots,M$) must be attributed to its best-supporting chunk $d^*_j \in \mathcal{R}(q)$ via:
\begin{equation}
d^*_j = \arg\max_{d_i \in \mathcal{R}(q)} \mathrm{sim}(\mathbf{e}_{s_j}, \mathbf{e}_i)
\label{eq:attribution}
\end{equation}
where $\mathbf{e}_{s_j}$ is the embedding of sentence $s_j$. Subsequently, a multi-signal hallucination detector—combining NLI entailment, BERTScore semantic similarity, and cosine overlap—evaluates whether $d^*_j$ supports $s_j$. Sentences failing verification are suppressed before reaching the user.

The computational bottleneck arises from three sequential sources:
\begin{enumerate}
    \item \textbf{Retrieval latency} $\mathcal{L}_{\mathrm{ret}}$: Scales as $O(Nd)$ for exhaustive search over $N$ chunks in $d$-dimensional space.
    \item \textbf{Generation latency} $\mathcal{L}_{\mathrm{gen}}$: LLM inference time (typically 3--7 seconds for 7B-parameter models under 4-bit quantization on CPU).
    \item \textbf{Verification latency} $\mathcal{L}_{\mathrm{NLI}}$: Cross-encoder NLI inference for $M$ generated sentences against $K$ retrieved chunks, scaling as $O(MK)$.
\end{enumerate}

Total per-query latency is:
\begin{equation}
\mathcal{L}_{\mathrm{total}} = \mathcal{L}_{\mathrm{ret}} + \mathcal{L}_{\mathrm{gen}} + M \cdot K \cdot \mathcal{L}_{\mathrm{NLI}}
\label{eq:total_latency}
\end{equation}

For a baseline system with $N{=}9{,}800$ chunks, $d{=}384$ (all-MiniLM-L6-v2 embeddings \cite{reimers2019sbert}), $K{=}5$, and $M{\approx}4$ sentences per answer, we observed $\mathcal{L}_{\mathrm{ret}} \approx 47$\,ms using FAISS FlatIP \cite{johnson2019faiss}, $\mathcal{L}_{\mathrm{gen}} \approx 5.2$\,s, and $\mathcal{L}_{\mathrm{NLI}} \approx 4.56$\,s, yielding $\mathcal{L}_{\mathrm{total}} \approx 9.8$\,s. While retrieval contributes only 0.5\% of total latency in this regime, two factors motivate optimization: (1)~retrieval latency scales linearly with $N$, becoming prohibitive for corpora with $N \gg 10^4$, and (2)~attribution computations (Eq.~\ref{eq:attribution}) require additional similarity searches, compounding retrieval overhead.

\subsection{Contributions}

This paper makes the following contributions:
\begin{itemize}
    \item A systematic benchmark of four vector indexing strategies—FAISS FlatIP, IVF-PQ, HNSW, and hybrid dense-sparse retrieval—tailored to explainability-driven offline RAG under resource constraints.
    \item Empirical Pareto analysis of the latency-accuracy-memory trade-space, demonstrating that HNSW achieves 12.4$\times$ retrieval speedup at 93\% Recall@5, hybrid retrieval attains 96\% Recall@5, and IVF-PQ enables edge deployment at 57\,MB memory footprint.
    \item Deployment guidelines mapping use-case requirements (real-time explainability, maximum accuracy, memory constraints, auditability) to optimal indexing configurations.
    \item Latency decomposition analysis revealing that NLI verification—not retrieval—constitutes the dominant bottleneck (46.5\% of end-to-end latency), informing future research directions.
\end{itemize}

\section{Related Work}

\subsection{Hallucination in Large Language Models}

Hallucination taxonomies distinguish \emph{factuality hallucination} (outputs contradicting world knowledge) from \emph{faithfulness hallucination} (outputs contradicting provided context) \cite{ji2023survey,huang2025survey}. Benchmarks such as TruthfulQA \cite{lin2022truthfulqa}, HaluEval \cite{li2023halueval}, RAGTruth \cite{niu2024ragtruth}, and FEVER \cite{thorne2018fever} provide labeled test sets for detection research. Probing studies demonstrate that LLMs exhibit poor calibration regarding their knowledge boundaries \cite{mallen2023trust}, and theoretical analysis suggests hallucination is fundamentally inevitable due to lossy compression of world knowledge into finite parameters \cite{xu2024inevitable}.

\subsection{Retrieval-Augmented Generation}

RAG \cite{lewis2020rag} combines dense retrieval with seq2seq generation, enabling models to condition outputs on retrieved evidence. Variants include REALM \cite{guu2020realm} with end-to-end retrieval-generation training, RETRO \cite{borgeaud2022retro} with chunked cross-attention, and Atlas \cite{izacard2022atlas} with joint index-model optimization. Recent extensions explore active retrieval \cite{jiang2023flare}, corrective RAG \cite{yan2024crag}, and domain-adapted pipelines \cite{krainovskikh2026domain}. However, these systems focus on retrieval quality as an input-side concern; none integrate mandatory output-side verification that retrieved evidence actually supports generated claims.

\subsection{Hallucination Detection Methods}

Detection approaches span internal consistency checks—SelfCheckGPT \cite{manakul2023selfcheckgpt} sampling multiple outputs and measuring agreement, semantic entropy \cite{farquhar2024semantic} quantifying uncertainty—and external verification methods. Fact-decomposition frameworks such as FActScore \cite{min2023factscore} and SAFE \cite{wei2024longform} break generations into atomic claims and verify each against a knowledge source. NLI-based consistency detectors including FactCC \cite{kryscinski2020factcc} and SummaC \cite{laban2022summac} classify entailment relationships between generated text and source documents. Automated evaluation frameworks like RAGAS \cite{es2024ragas} and G-Eval \cite{liu2023geval} employ LLMs as judges. A common limitation is reliance on cloud APIs or heavyweight model calls, precluding offline deployment.

\subsection{Explainable AI for RAG}

XAI research for NLP \cite{minaee2024xai,rahman2026survey} has explored attention visualization, gradient-based attribution, and probing techniques. Recent work on RAG explainability includes FEEL \cite{zhang2024feel}, which generates post-hoc natural language explanations of retrieval-generation chains, and attributed generation frameworks \cite{agrawal2024attributed} that annotate outputs with source citations. Unlike prior work treating explainability as an analytical overlay, our framework makes attribution \emph{constitutive}: computed inline and gating what reaches the user.

\subsection{Approximate Nearest Neighbor Search}

Scaling dense retrieval to large corpora necessitates approximate nearest neighbor (ANN) algorithms. FAISS \cite{johnson2019faiss} provides optimized implementations of flat indexes (exact search), inverted file indexes (IVF) for coarse quantization, product quantization (PQ) for memory compression, and composite methods (IVF-PQ). Hierarchical Navigable Small World (HNSW) graphs \cite{malkov2018hnsw} achieve sub-millisecond retrieval via greedy graph traversal, demonstrating state-of-the-art accuracy-speed trade-offs on ANN benchmarks. Hybrid retrieval combining dense embeddings with sparse keyword matching (e.g., BM25 \cite{robertson2009bm25}) has been explored in ColBERTv2 \cite{santhanam2022colbertv2} and SPLADE \cite{formal2021splade}. Our contribution is the first Pareto analysis of these methods specifically for explainability-constrained offline RAG under resource limitations.

\section{System Architecture}

\subsection{Pipeline Overview}

The explainability-driven RAG pipeline comprises six stages executed sequentially for each query:

\textbf{Stage 1: Document Ingestion.} Source documents (PDF, DOCX, TXT) are parsed with page-boundary preservation and segmented into overlapping chunks using a 512-token sliding window with 50-token overlap, balancing sentence encoder context limits against evidentiary coherence.

\textbf{Stage 2: Dense Retrieval.} The query embedding $\mathbf{e}_q$ is computed via all-MiniLM-L6-v2 (384-dimensional, $\ell_2$-normalized) and searched against the vector index to retrieve Top-$K{=}5$ chunks $\mathcal{R}(q)$ with highest cosine similarity.

\textbf{Stage 3: XAI Evidence Attribution Layer.} For each retrieved chunk $d_i \in \mathcal{R}(q)$, attribution scores $\alpha_i = \mathrm{sim}(\mathbf{e}_q, \mathbf{e}_i)$ are computed and stored. Chunk provenance metadata (source file, page number) is preserved for citation generation.

\textbf{Stage 4: Local LLM Generation.} A locally hosted model (Mistral-7B-Instruct or LLaMA-3-8B, 4-bit quantized via Ollama) generates a draft answer conditioned on the concatenation of retrieved chunks, using temperature 0.1 and maximum 512 tokens to reduce stochastic extrapolation.

\textbf{Stage 5: Multi-Signal Hallucination Detection.} The draft is split into sentences $\{s_1, \ldots, s_M\}$. For each $s_j$:
\begin{itemize}
    \item Best-supporting chunk $d^*_j$ is identified via Eq.~\ref{eq:attribution}.
    \item Three detectors score the $(d^*_j, s_j)$ pair:
    \begin{itemize}
        \item NLI entailment: $E_j = P(\mathrm{entailment} \mid d^*_j, s_j)$ via cross-encoder/nli-deberta-v3-small.
        \item BERTScore: $B_j = \mathrm{F1}(d^*_j, s_j)$ using DistilBERT backbone.
        \item Cosine similarity: $C_j = \mathrm{sim}(\mathbf{e}_{s_j}, \mathbf{e}_{d^*_j})$.
    \end{itemize}
    \item Scores are fused: $S_j = 0.45 E_j + 0.35 B_j + 0.20 C_j$.
\end{itemize}

\textbf{Stage 6: Tri-State Decision Controller.} Each sentence is classified based on $S_j$:
\begin{equation}
D(s_j) =
\begin{cases}
\textsc{Supported}, & S_j \geq 0.75 \\
\textsc{Weakly\_Supported}, & 0.50 \leq S_j < 0.75 \\
\textsc{Hallucinated}, & S_j < 0.50
\end{cases}
\label{eq:tristate}
\end{equation}
\textsc{Supported} sentences are emitted with inline citations; \textsc{Weakly\_Supported} carry warning flags; \textsc{Hallucinated} are suppressed. The entire answer is blocked if $>$50\% of sentences are hallucinated.

This architecture makes hallucination detection \emph{intrinsic} rather than post hoc: no generated text reaches the user without passing verification.

\subsection{Baseline Retrieval: FAISS FlatIP}

FAISS FlatIP performs exhaustive $\ell_2$-normalized inner product search, mathematically equivalent to cosine similarity for unit vectors:
\begin{equation}
\mathrm{FlatIP}(\mathbf{e}_q) = \arg\max_{i=1,\ldots,N} \mathbf{e}_q \cdot \mathbf{e}_i, \quad \|\mathbf{e}_q\| = \|\mathbf{e}_i\| = 1
\label{eq:flatip}
\end{equation}

Computational complexity is $O(Nd)$ per query with zero approximation error. Memory footprint is $4Nd$ bytes (float32 storage). For $N{=}9{,}800$ chunks and $d{=}384$, this yields a 15.05\,MB index. FlatIP provides deterministic, auditable retrieval—critical for explainability applications requiring reproducible provenance traces. However, latency scales linearly with corpus size, motivating approximate methods.

\section{Optimization Methodology}

\subsection{Inverted File Index with Product Quantization (IVF-PQ)}

\subsubsection{Inverted File Structure}
IVF partitions the embedding space into $n_{\mathrm{list}}$ Voronoi cells via $k$-means clustering. Each cell is represented by a centroid $\mathbf{c}_\ell$, and vectors are assigned to their nearest centroid:
\begin{equation}
\mathrm{assign}(\mathbf{e}_i) = \arg\min_{\ell=1,\ldots,n_{\mathrm{list}}} \|\mathbf{e}_i - \mathbf{c}_\ell\|_2
\label{eq:ivf_assign}
\end{equation}

At query time, only the $n_{\mathrm{probe}}$ nearest lists (by centroid distance) are searched, reducing complexity to:
\begin{equation}
O\left(\frac{N}{n_{\mathrm{list}}} \cdot n_{\mathrm{probe}} \cdot d\right)
\label{eq:ivf_complexity}
\end{equation}

This coarse quantization trades recall for speed: vectors in un-probed lists cannot be retrieved, introducing false negatives. Increasing $n_{\mathrm{probe}}$ recovers recall at latency cost.

\subsubsection{Product Quantization}
Each residual vector $\mathbf{r}_i = \mathbf{e}_i - \mathbf{c}_{\mathrm{assign}(\mathbf{e}_i)}$ is divided into $m$ subvectors of dimension $d/m$ and quantized independently using learned codebooks $\mathcal{C}_1, \ldots, \mathcal{C}_m$, each containing $2^b$ centroids:
\begin{equation}
\mathrm{PQ}(\mathbf{r}_i) = [\mathrm{code}_1, \ldots, \mathrm{code}_m], \quad \mathrm{code}_j \in \{0, \ldots, 2^b{-}1\}
\label{eq:pq}
\end{equation}

The quantized vector is reconstructed as:
\begin{equation}
\hat{\mathbf{r}}_i = [\mathcal{C}_1[\mathrm{code}_1], \ldots, \mathcal{C}_m[\mathrm{code}_m]]
\label{eq:pq_reconstruct}
\end{equation}

Expected reconstruction error is:
\begin{equation}
\mathbb{E}[\|\mathbf{r}_i - \hat{\mathbf{r}}_i\|_2^2] = \sum_{j=1}^{m} \mathbb{E}[\|\mathbf{r}_i^{(j)} - \mathcal{C}_j[\mathrm{code}_j]\|_2^2]
\label{eq:pq_error}
\end{equation}

Memory footprint reduces to $mb$ bits per vector. For $m{=}8$ subspaces and $b{=}8$ bits, storage is 8 bytes per vector—96$\times$ compression compared to 384 bytes (float32). Distance computations operate on precomputed lookup tables, enabling fast approximate similarity scoring.

\subsubsection{Configuration and Trade-offs}
We configure $n_{\mathrm{list}}{=}128$ (empirically effective for $N{\sim}10^4$), $n_{\mathrm{probe}}{=}16$ (2--8$\times$ speedup vs. exhaustive with $<$10\% recall loss), $m{=}8$ subspaces (384 dims $\div$ 8 = 48 dims/subspace), and $b{=}8$ bits (256 centroids per codebook). This balances memory compression, query speed, and recall degradation.

\subsection{Hierarchical Navigable Small World (HNSW)}

HNSW constructs a multi-layer proximity graph where nodes represent embeddings and edges connect approximate nearest neighbors. The graph has $L$ layers ($L \approx \log_2 N$), with layer 0 containing all vectors and each higher layer containing an exponentially decreasing subset.

Search proceeds via greedy best-first traversal:
\begin{equation}
\mathrm{HNSW}(\mathbf{e}_q) : \text{greedy-search on layers } \ell = L, \ldots, 0
\label{eq:hnsw_search}
\end{equation}

Starting from a random entry point at the top layer, the algorithm iteratively moves to the neighbor closest to $\mathbf{e}_q$ until a local minimum is reached, then descends to the next layer. At layer 0, a beam search with width $ef_{\mathrm{search}}$ explores the $K$ nearest neighbors.

Construction complexity is $O(N \log N \cdot M \cdot d)$, where $M$ is the maximum number of edges per node. Query complexity is $O(\log N \cdot ef_{\mathrm{search}} \cdot d)$, achieving sub-linear scaling. HNSW demonstrates state-of-the-art performance on ANN-Benchmarks \cite{aumuller2020annbench}, with sub-millisecond retrieval for $N \sim 10^6$ on CPU.

Key parameters:
\begin{itemize}
    \item $M$: Max bidirectional links per node (higher $M$ improves recall but increases memory and build time). We use $M{=}32$.
    \item $ef_{\mathrm{construction}}$: Beam width during index building (higher yields better graph connectivity). We use $ef_{\mathrm{construction}}{=}200$.
    \item $ef_{\mathrm{search}}$: Beam width during query (primary latency-recall tuning knob). We sweep $ef_{\mathrm{search}} \in \{16, 32, 64, 128\}$ to trace the Pareto frontier.
\end{itemize}

\subsection{Hybrid Dense-Sparse Retrieval}

Dense embeddings capture semantic similarity but may miss exact keyword matches critical for factual verification (e.g., technical terminology, proper nouns, command syntax). We augment FAISS dense retrieval with BM25 sparse retrieval \cite{robertson2009bm25}.

BM25 scores document $d_i$ for query $q$ as:
\begin{equation}
\mathrm{BM25}(q, d_i) = \sum_{t \in q} \mathrm{IDF}(t) \cdot \frac{f(t, d_i) \cdot (k_1 + 1)}{f(t, d_i) + k_1 \cdot \left(1 - b + b \cdot \frac{|d_i|}{\mathrm{avgdl}}\right)}
\label{eq:bm25}
\end{equation}
where:
\begin{itemize}
    \item $f(t, d_i)$: term frequency of token $t$ in document $d_i$
    \item $\mathrm{IDF}(t) = \log\frac{N - n(t) + 0.5}{n(t) + 0.5}$: inverse document frequency
    \item $|d_i|$: document length; $\mathrm{avgdl}$: average document length
    \item $k_1{=}1.5$, $b{=}0.75$: tuning constants (standard values)
\end{itemize}

Dense and sparse scores are fused via weighted combination:
\begin{equation}
\mathrm{score}_{\mathrm{hybrid}}(q, d_i) = \alpha \cdot \mathrm{sim}(\mathbf{e}_q, \mathbf{e}_i) + (1-\alpha) \cdot \frac{\mathrm{BM25}(q, d_i)}{\max_{j} \mathrm{BM25}(q, d_j)}
\label{eq:hybrid}
\end{equation}

We normalize BM25 scores to $[0,1]$ by dividing by the per-query maximum and set $\alpha{=}0.7$ (dense-weighted) based on ablation experiments on a 50-query development set. BM25 retrieval runs on a separate inverted index (implementation: Elasticsearch 7.17), adding $\mathcal{L}_{\mathrm{BM25}} \approx 2$--5\,ms but improving Recall@5 by 2--4 percentage points for keyword-heavy technical queries.

\section{Experimental Setup}

\subsection{Corpus and Indexing}

\textbf{Dataset:} We curate a corpus of 120 technical reference documents—Git version control guides, Docker container documentation, and Python developer manuals—totaling approximately 6.7\,MB of text. Documents are segmented into $N{=}9{,}800$ chunks using a 512-token sliding window with 50-token overlap. Each chunk is embedded with \texttt{sentence-transformers/all-MiniLM-L6-v2} (22.7M parameters, 384-dimensional output), yielding $\ell_2$-normalized vectors suitable for inner-product search.

\textbf{Indexes constructed:}
\begin{itemize}
    \item \textbf{FlatIP (baseline):} FAISS IndexFlatIP, exact exhaustive search.
    \item \textbf{IVF-PQ:} FAISS IndexIVFPQ with $n_{\mathrm{list}}{=}128$, $n_{\mathrm{probe}}{=}16$, $m{=}8$ subspaces, $b{=}8$ bits per code.
    \item \textbf{HNSW:} FAISS IndexHNSWFlat with $M{=}32$, $ef_{\mathrm{construction}}{=}200$, $ef_{\mathrm{search}} \in \{16, 32, 64, 128\}$ (swept at query time).
    \item \textbf{Hybrid:} FAISS IndexFlatIP + Elasticsearch BM25 inverted index, fusion weight $\alpha{=}0.7$.
\end{itemize}

Indexes are built offline once and serialized to disk. Build times: FlatIP 0.8\,s, IVF-PQ 14.3\,s (includes $k$-means clustering and PQ codebook training), HNSW 18.2\,min, Elasticsearch indexing 3.1\,s.

\subsection{Evaluation Metrics}

\subsubsection{Retrieval Quality}
We measure Recall@$K$, defined as:
\begin{equation}
\mathrm{Recall@}K = \frac{|\mathcal{R}_{\mathrm{approx}}(q) \cap \mathcal{R}_{\mathrm{exact}}(q)|}{K}
\label{eq:recall}
\end{equation}
where $\mathcal{R}_{\mathrm{exact}}(q)$ is the FlatIP Top-$K$ result (ground truth) and $\mathcal{R}_{\mathrm{approx}}(q)$ is the approximate index result. Averaged over 400 test queries comprising 250 answerable questions (evidence present in corpus) and 150 adversarial queries (out-of-domain, premise-corrupted, or unanswerable) to stress-test robustness.

\subsubsection{Latency}
Mean and 95th-percentile retrieval time (milliseconds) measured via Python \texttt{timeit} module with 1000 warm-up iterations followed by 500 timed runs per query. Embedding computation time is excluded (embeddings cached) to isolate index search overhead. Reported latencies include only the vector search operation (FAISS index query or Elasticsearch BM25 scoring).

\subsubsection{Memory}
Index file size on disk (megabytes) and resident RAM consumption during queries (measured via \texttt{psutil.Process().memory\_info().rss}).

\subsubsection{End-to-End Impact}
Total query latency $\mathcal{L}_{\mathrm{total}}$ (Eq.~\ref{eq:total_latency}) including:
\begin{itemize}
    \item Retrieval: FAISS + optional BM25
    \item Generation: Mistral-7B-Instruct-v0.2 (4-bit GGUF) via Ollama
    \item Verification: cross-encoder/nli-deberta-v3-small (DeBERTa-v3, 140M parameters) for $M{\approx}4$ sentences $\times$ $K{=}5$ chunks = 20 cross-encodings per query
\end{itemize}

\subsection{Hardware Platform}

All experiments conducted on a single machine:
\begin{itemize}
    \item \textbf{CPU:} Intel Core i7-11700 (8 cores, 16 threads, 2.5\,GHz base / 4.9\,GHz turbo)
    \item \textbf{RAM:} 16\,GB DDR4-3200
    \item \textbf{GPU:} None (CPU-only to simulate resource-constrained deployment)
    \item \textbf{OS:} Ubuntu 22.04 LTS, Linux kernel 5.15
    \item \textbf{Software:} Python 3.9, FAISS 1.7.4 (CPU build), PyTorch 2.0.1, Transformers 4.30, Ollama 0.1.29
\end{itemize}

All FAISS operations restricted to single-threaded execution (\texttt{OMP\_NUM\_THREADS=1}) to isolate algorithmic efficiency and ensure reproducibility. NLI cross-encoder and LLM generation use multi-threading as permitted by default PyTorch/Ollama configurations.

\section{Results and Analysis}

\subsection{Retrieval Latency and Accuracy}

Table~\ref{tab:results} summarizes retrieval performance across all configurations. HNSW with $ef_{\mathrm{search}}{=}64$ achieves the optimal latency-recall balance: 3.8\,ms mean latency (12.4$\times$ speedup relative to FlatIP baseline) while maintaining Recall@5 = 0.93. IVF-PQ attains 7.1\,ms (6.7$\times$ speedup) but Recall@5 degrades to 0.89, marginally below the 0.90 threshold recommended for production RAG systems \cite{karpukhin2020dense}. The hybrid dense-sparse approach reaches Recall@5 = 0.96 (highest among all methods) at 8.2\,ms, demonstrating that sparse keyword matching recovers chunks missed by semantic embeddings—particularly for technical queries containing domain-specific terminology.

95th-percentile latencies follow similar trends, with HNSW ($ef{=}64$) at 5.9\,ms and hybrid at 11.3\,ms, indicating stable performance under tail latency. Memory footprints reveal a three-way trade-off: IVF-PQ achieves 6.7$\times$ compression (57\,MB vs. 384\,MB baseline), HNSW increases memory to 412\,MB due to graph storage overhead, and hybrid requires 421\,MB (FAISS + Elasticsearch inverted index).

\begin{table}[t]
\caption{Retrieval Performance Comparison ($N{=}9{,}800$, $K{=}5$, $d{=}384$)}
\label{tab:results}
\centering
\small
\begin{tabular}{@{}lccccc@{}}
\toprule
Method & \begin{tabular}[c]{@{}c@{}}Latency\\(ms)\end{tabular} & \begin{tabular}[c]{@{}c@{}}95\%ile\\(ms)\end{tabular} & Recall@5 & \begin{tabular}[c]{@{}c@{}}Memory\\(MB)\end{tabular} & Speedup \\
\midrule
FlatIP (baseline) & 47.3 & 52.1 & 1.000 & 384 & 1.0$\times$ \\
\midrule
IVF-PQ & 7.1 & 9.8 & 0.89 & 57 & 6.7$\times$ \\
HNSW ($ef{=}16$) & 2.1 & 3.2 & 0.87 & 412 & 22.5$\times$ \\
HNSW ($ef{=}32$) & 3.1 & 4.5 & 0.91 & 412 & 15.3$\times$ \\
HNSW ($ef{=}64$) & \textbf{3.8} & \textbf{5.9} & \textbf{0.93} & 412 & 12.4$\times$ \\
HNSW ($ef{=}128$) & 6.4 & 9.1 & 0.95 & 412 & 7.4$\times$ \\
\midrule
Hybrid (FAISS+BM25) & 8.2 & 11.3 & \textbf{0.96} & 421 & 5.8$\times$ \\
\bottomrule
\end{tabular}
\end{table}

\subsection{Latency-Accuracy Pareto Frontier}

Plotting latency versus Recall@5 reveals that HNSW configurations dominate the Pareto frontier (Fig.~\ref{fig:pareto_placeholder}). The $ef{=}64$ operating point represents the optimal balance for real-time explainability applications: sub-5\,ms retrieval with $>$0.90 recall. Lower $ef$ values (16, 32) sacrifice recall for speed, acceptable only if downstream verification is highly robust. Higher $ef$ (128) approaches FlatIP recall (0.95) but at diminishing latency benefits.

IVF-PQ lies off the Pareto curve—HNSW ($ef{=}32$) achieves comparable recall (0.91 vs. 0.89) at lower latency (3.1\,ms vs. 7.1\,ms). However, IVF-PQ's 6.7$\times$ memory compression makes it the sole viable option for memory-constrained edge devices ($<$100\,MB RAM budget).

Hybrid retrieval sits at the high-accuracy extreme: 0.96 recall justifies the 8.2\,ms latency for applications prioritizing correctness over speed, such as batch processing or high-stakes document review.

\subsection{End-to-End System Latency}

Table~\ref{tab:e2e} decomposes total query latency into retrieval, generation, and NLI verification components. Baseline FlatIP yields $\mathcal{L}_{\mathrm{total}}{=}9.81$\,s (47.3\,ms retrieval + 5.20\,s generation + 4.56\,s NLI). Optimizing retrieval via HNSW ($ef{=}64$) reduces total latency to 9.76\,s—a 0.5\% improvement. Hybrid achieves 9.77\,s.

\begin{table}[t]
\caption{End-to-End Query Latency Breakdown (mean, 400 queries)}
\label{tab:e2e}
\centering
\small
\begin{tabular}{@{}lcccc@{}}
\toprule
Configuration & $\mathcal{L}_{\mathrm{ret}}$ (s) & $\mathcal{L}_{\mathrm{gen}}$ (s) & $\mathcal{L}_{\mathrm{NLI}}$ (s) & Total (s) \\
\midrule
FlatIP & 0.047 & 5.20 & 4.56 & 9.81 \\
HNSW ($ef{=}64$) & 0.004 & 5.20 & 4.56 & 9.76 \\
Hybrid & 0.008 & 5.20 & 4.56 & 9.77 \\
\bottomrule
\end{tabular}
\end{table}

This modest end-to-end improvement—despite 12.4$\times$ retrieval speedup—reveals that \textbf{NLI verification is the dominant bottleneck}, consuming 46.5\% of total latency. LLM generation accounts for 53.1\%, while retrieval contributes only 0.5\% in the baseline and becomes negligible (0.04\%) with HNSW. Consequently, further end-to-end acceleration requires addressing NLI overhead via model distillation, early-exit strategies, or bi-encoder approximations—beyond the scope of this retrieval-focused study but critical for future work.

\subsection{Recall@K Sensitivity Analysis}

We evaluate Recall@$K$ for $K \in \{1, 3, 5, 10, 20\}$ to assess robustness across retrieval set sizes (Fig.~\ref{fig:recall_k_placeholder}). HNSW ($ef{=}64$) maintains $\geq$0.92 recall for all $K$, indicating consistent top-rank precision. IVF-PQ degrades more sharply: Recall@10 drops to 0.83 and Recall@20 to 0.78, confirming that coarse quantization introduces false negatives that compound at larger $K$. This has implications for multi-hop reasoning systems requiring broad evidence retrieval.

Hybrid retrieval consistently outperforms by 2--3 percentage points across all $K$, validating the hypothesis that dense-sparse fusion recovers complementary failure modes. At $K{=}1$ (single best chunk), hybrid achieves 0.89 versus 0.85 for HNSW and 0.81 for IVF-PQ—critical for citation-generation tasks requiring precise provenance.

\subsection{Memory-Constrained Deployment}

For embedded systems with $<$100\,MB RAM budget (e.g., Raspberry Pi 4 with 4\,GB total RAM allocated across OS, application, and models), IVF-PQ is the only viable indexing strategy. The 57\,MB footprint leaves headroom for embedder and lightweight models.

While baseline IVF-PQ ($n_{\mathrm{probe}}{=}16$) yields Recall@5 = 0.89, increasing $n_{\mathrm{probe}}{=}24$ recovers Recall@5 = 0.91 at 9.3\,ms latency—still 5.1$\times$ faster than FlatIP and meeting the 0.90 production threshold. This configuration enables offline explainability-driven RAG on edge devices, critical for IoT, mobile, or air-gapped deployments.

\section{Discussion}

\subsection{Deployment Guidelines}

Based on empirical Pareto analysis, we provide the following decision tree for practitioners:

\begin{itemize}
    \item \textbf{Real-time explainability ($<$5\,ms retrieval latency):} HNSW with $ef_{\mathrm{search}}{=}64$ (Recall@5 = 0.93, 3.8\,ms). Suitable for interactive systems requiring per-sentence provenance tracing.
    
    \item \textbf{Maximum accuracy (batch processing acceptable):} Hybrid FAISS+BM25 (Recall@5 = 0.96, 8.2\,ms). Optimal for high-stakes document review, legal analysis, or compliance verification where precision outweighs latency.
    
    \item \textbf{Memory-constrained edge deployment ($<$100\,MB RAM):} IVF-PQ with $n_{\mathrm{probe}}{=}24$ (Recall@5 = 0.91, 9.3\,ms, 57\,MB). Enables offline RAG on Raspberry Pi, Jetson Nano, or mobile devices.
    
    \item \textbf{Auditable systems (regulatory compliance, forensics):} FlatIP (Recall@5 = 1.00, deterministic). Eliminates approximation artifacts; reproducible provenance for legal admissibility.
\end{itemize}

\subsection{Limitations and Future Work}

\textbf{HNSW index maintenance:} Construction time (18.2 minutes for 9,800 vectors) precludes real-time corpus updates. Incremental insertion is supported but degrades graph connectivity over time, requiring periodic rebuilds. For dynamic corpora, IVF-PQ with periodic re-clustering may be preferable.

\textbf{IVF-PQ domain transfer:} Product quantization codebooks are corpus-dependent. Codebooks trained on technical documentation may not transfer to other domains (e.g., medical literature, legal contracts) without retraining, introducing deployment friction.

\textbf{Hybrid retrieval complexity:} Maintaining parallel FAISS and BM25 indexes doubles storage overhead and complicates corpus ingestion pipelines. Fusion weight $\alpha$ may require per-domain tuning; we used a single $\alpha{=}0.7$ across all queries.

\textbf{NLI bottleneck:} While retrieval optimization yields 12.4$\times$ speedup, end-to-end latency improves by only 0.5\% due to NLI verification overhead (4.56\,s per query). Future work must address this via:
\begin{itemize}
    \item \textbf{Bi-encoder approximation:} Replace cross-encoder NLI with faster bi-encoder similarity, trading precision for 10--50$\times$ speedup.
    \item \textbf{Early-exit mechanisms:} Abort NLI inference for sentences with high cosine/BERTScore (likely supported).
    \item \textbf{Model distillation:} Compress DeBERTa-v3 (140M params) to DistilBERT-scale (66M) while preserving entailment accuracy.
\end{itemize}

Targeting $\mathcal{L}_{\mathrm{NLI}} < 1$\,s would reduce $\mathcal{L}_{\mathrm{total}}$ to $\sim$6\,s, achieving practical real-time explainability.

\subsection{Implications for Explainability-Driven RAG}

Retrieval latency optimization directly impacts explainability usability: reducing $\mathcal{L}_{\mathrm{ret}}$ from 47\,ms to 4\,ms per query enables real-time per-sentence provenance tracing in interactive systems. However, the Pareto analysis reveals a broader architectural insight: \textbf{retrieval is not the bottleneck in explainability-constrained pipelines}—verification is. This inverts conventional RAG optimization priorities, which focus on scaling retrieval to millions of documents \cite{izacard2022atlas} while treating hallucination detection as a lightweight post-processing step.

Our findings suggest that future explainability-augmented RAG systems should co-optimize retrieval and verification latencies rather than treating them independently. Hybrid architectures combining fast approximate retrieval (HNSW) with approximate verification (bi-encoder NLI) may achieve the latency-accuracy balance required for production deployment.

\section{Conclusion}

This paper presented a comprehensive evaluation of vector indexing strategies for explainability-driven offline RAG systems under resource constraints. Through systematic benchmarking of FAISS FlatIP, IVF-PQ, HNSW, and hybrid dense-sparse retrieval across latency, memory, and Recall@K metrics, we identified optimal operating points for diverse deployment scenarios:

\begin{itemize}
    \item HNSW with tuned $ef_{\mathrm{search}}{=}64$ achieves 12.4$\times$ retrieval speedup at 93\% recall, enabling sub-5\,ms provenance attribution for real-time explainability.
    \item Hybrid FAISS+BM25 attains 96\% recall—the highest among tested methods—demonstrating that dense-sparse fusion recovers complementary retrieval failures.
    \item IVF-PQ enables edge deployment at 57\,MB memory footprint with acceptable 91\% recall, supporting offline RAG on resource-constrained devices.
\end{itemize}

Critically, end-to-end latency decomposition revealed that retrieval optimization yields only 0.5\% total speedup due to the dominant NLI verification bottleneck (46.5\% of latency). This finding redirects future research toward co-optimizing retrieval and verification, potentially via bi-encoder approximations, early-exit strategies, or model distillation targeting sub-3\,s end-to-end query latency for practical real-time explainability.

Our Pareto analysis and deployment guidelines provide actionable insights for practitioners balancing latency, accuracy, and memory in production explainability-augmented RAG systems, while identifying NLI acceleration as the critical path for next-generation offline hallucination detection frameworks.

\bibliographystyle{IEEEtran}
\bibliography{references}

\end{document}
